% Results margin figure: each circuit's charging and discharging tau from
% the scope fits on one axis, so the gap between them reads as a length.
% Charging sits just below its row's bar and discharging just above it, both
% pointing at it; where the two agree they meet tip to tip. Each row's tags
% give the gap in seconds (on the head line) and percent (one line under it),
% set as a single two-line block per row, all on one decimal column.
%
% Data: data/scope-values.tex (fitted tau over the nominal R*C), written by
% make-scope-data.py. x is (tau / R*C - 1) in percent:
%   Circuit 1  0.9566 -> -4.3 %   1.0592 -> +5.9 %   1.0592/0.9566 = 1.107
%   Circuit 2  0.9633 -> -3.7 %   0.9679 -> -3.2 %   0.9679/0.9633 = 1.005
%   Circuit 3  0.9783 -> -2.2 %   0.9784 -> -2.2 %   0.9784/0.9783 = 1.000
%   Circuit 4  0.9984 -> -0.2 %   1.0334 -> +3.3 %   1.0334/0.9984 = 1.035
% The seconds are the same ratios times R*C (45.0/49.8, 96.3/96.8,
% 143.8/143.8, 31.9/33.0 s), as in the Part 4 tables, and the gaps
% 49.78 - 44.96 = 4.82, 96.79 - 96.33 = 0.46, 143.82 - 143.81 = 0.01 and
% 33.04 - 31.92 = 1.12 s. Adjusted tau:
% (R || R_VM)/R = 9.598/10.598 = 0.9056, so -9.4 %.
% Written by make-scope-data.py; do not edit.
% Zero offset 0.408 V. tau is the fitted tau over the nominal R*C;
% level is 63.2 % of the way from the plateau to the fitted final value.
% c1: 47s-c_1.csv
\expandafter\def\csname scope-c1-tau\endcsname{0.9566}
\expandafter\def\csname scope-c1-taus\endcsname{45.0}
\expandafter\def\csname scope-c1-err\endcsname{-4.3}
\expandafter\def\csname scope-c1-level\endcsname{8.280}
\expandafter\def\csname scope-c1-vinf\endcsname{13.09}
\expandafter\def\csname scope-c1-vzero\endcsname{0.02}
% c2: 100s-c.csv
\expandafter\def\csname scope-c2-tau\endcsname{0.9633}
\expandafter\def\csname scope-c2-taus\endcsname{96.3}
\expandafter\def\csname scope-c2-err\endcsname{-3.7}
\expandafter\def\csname scope-c2-level\endcsname{8.346}
\expandafter\def\csname scope-c2-vinf\endcsname{13.22}
\expandafter\def\csname scope-c2-vzero\endcsname{-0.03}
% c3: 147s-c-0002_1.csv
\expandafter\def\csname scope-c3-tau\endcsname{0.9783}
\expandafter\def\csname scope-c3-taus\endcsname{144}
\expandafter\def\csname scope-c3-err\endcsname{-2.2}
\expandafter\def\csname scope-c3-level\endcsname{8.417}
\expandafter\def\csname scope-c3-vinf\endcsname{13.32}
\expandafter\def\csname scope-c3-vzero\endcsname{-0.00}
% c4: 32s-c.csv
\expandafter\def\csname scope-c4-tau\endcsname{0.9984}
\expandafter\def\csname scope-c4-taus\endcsname{31.9}
\expandafter\def\csname scope-c4-err\endcsname{-0.2}
\expandafter\def\csname scope-c4-level\endcsname{8.547}
\expandafter\def\csname scope-c4-vinf\endcsname{13.51}
\expandafter\def\csname scope-c4-vzero\endcsname{0.02}
% d1: 47s-d.csv
\expandafter\def\csname scope-d1-tau\endcsname{1.0592}
\expandafter\def\csname scope-d1-taus\endcsname{49.8}
\expandafter\def\csname scope-d1-err\endcsname{+5.9}
\expandafter\def\csname scope-d1-level\endcsname{5.502}
\expandafter\def\csname scope-d1-vinf\endcsname{0.02}
\expandafter\def\csname scope-d1-vzero\endcsname{14.92}
% d2: 100s-d_1.csv
\expandafter\def\csname scope-d2-tau\endcsname{0.9679}
\expandafter\def\csname scope-d2-taus\endcsname{96.8}
\expandafter\def\csname scope-d2-err\endcsname{-3.2}
\expandafter\def\csname scope-d2-level\endcsname{5.453}
\expandafter\def\csname scope-d2-vinf\endcsname{-0.04}
\expandafter\def\csname scope-d2-vzero\endcsname{14.88}
% d3: 143s-d.csv
\expandafter\def\csname scope-d3-tau\endcsname{0.9784}
\expandafter\def\csname scope-d3-taus\endcsname{144}
\expandafter\def\csname scope-d3-err\endcsname{-2.2}
\expandafter\def\csname scope-d3-level\endcsname{5.515}
\expandafter\def\csname scope-d3-vinf\endcsname{0.06}
\expandafter\def\csname scope-d3-vzero\endcsname{14.88}
% d4: 32s-d.csv
\expandafter\def\csname scope-d4-tau\endcsname{1.0334}
\expandafter\def\csname scope-d4-taus\endcsname{33.0}
\expandafter\def\csname scope-d4-err\endcsname{+3.3}
\expandafter\def\csname scope-d4-level\endcsname{5.452}
\expandafter\def\csname scope-d4-vinf\endcsname{-0.05}
\expandafter\def\csname scope-d4-vzero\endcsname{14.90}

\newcommand{\prtau}[2]{\csname scope-#1#2-tau\endcsname}   % c|d, circuit
\newcommand{\prx}[2]{{(\prtau{#1}{#2}-1)*100}}
% The marks as text, for the legend and row heads
\newcommand{\chgmark}{\tikz[baseline=0pt]\fill[traceInk]
  (0,0) -- (5pt,0) -- (2.5pt,4.5pt) -- cycle;}
\newcommand{\dismark}{\tikz[baseline=0pt]\fill[traceInk]
  (0,4.5pt) -- (5pt,4.5pt) -- (2.5pt,0) -- cycle;}
% Fitted tau in seconds: ratio times R*C
\newcommand{\prsec}[3]{\pgfmathparse{\prtau{#1}{#2}*#3}%
  \pgfmathprintnumber[fixed, fixed zerofill, precision=1,
    assume math mode=true]{\pgfmathresult}\,\unit{\second}}
% How much longer the discharge took than the charge, signed, to one
% decimal; under 0.05 prints as "< 0.1". Set as two nodes split at the
% decimal point, which sits \prdecx left of the row's right edge, so every
% tag lines up on it. \prdelta{expression}{unit}{y shift}{circuit}
\newlength{\prdecx}
\newcommand{\prdelta}[4]{\pgfmathsetmacro{\prval}{#1}%
  \pgfmathtruncatemacro{\prtenths}{round(\prval*10)}%
  \ifnum\prtenths<1 \def\prwhole{<\,0}\def\prfrac{1}%
  \else
    \pgfmathtruncatemacro{\prwhole}{\prtenths/10}%
    \pgfmathtruncatemacro{\prfrac}{\prtenths-10*\prwhole}%
    \edef\prwhole{+\prwhole}%
  \fi
  \node[prtag, inner sep=0pt, anchor=base east]
    at ([xshift=\prdecx, yshift=#3]prowr#4) {\prwhole};
  \node[prtag, inner sep=0pt, anchor=base west]
    at ([xshift=\prdecx, yshift=#3]prowr#4) {.\prfrac\,\unit{#2}};}
\newcommand{\prgap}[2]{\prdelta{(\prtau{d}{#1}/\prtau{c}{#1}-1)*100}{\percent}{#2}{#1}}
\newcommand{\prdiff}[3]{\prdelta{(\prtau{d}{#1}-\prtau{c}{#1})*#2}{\second}{#3}{#1}}
% One row: the gap bar on the line, and its ends saved for the marks,
% which are drawn after the axis (pt offsets inside it go astray).
\newcommand{\pairrow}[2]{% circuit, row y
  \draw[accent, line width=1.6pt] (axis cs:\prx{c}{#1},#2) -- (axis cs:\prx{d}{#1},#2);
  \coordinate (prc#1) at (axis cs:\prx{c}{#1},#2);
  \coordinate (prd#1) at (axis cs:\prx{d}{#1},#2);
  \coordinate (prow#1) at (axis cs:-12,#2);
  \coordinate (prowr#1) at (axis cs:12,#2);}
\tikzset{
  prhead/.style={note, anchor=base west, fill=white, inner sep=1pt},
  prtag/.style={tag, text=chartText, anchor=base east, fill=white, inner sep=1pt},
}

% Decimal column: the widest fraction (".0 %") ends 1 pt in from the edge.
% Measured outside the picture (TikZ sets \nullfont inside one), with the
% chart's siunitx settings so the unit is in its Myriad.
\settowidth{\prdecx}{\sisetup{mode=text, reset-text-family=false,
  reset-text-series=false}\sffamily\bfseries\fontsize{7.5}{9}\selectfont
  .0\,\unit{\percent}}
\setlength{\prdecx}{\dimexpr-\prdecx-1pt\relax}

\begin{tikzpicture}
  \begin{axis}[labplot, name=pairs, width=158pt, height=118pt,
      xmin=-12, xmax=12, ymin=0.55, ymax=4.75,
      axis x line*=bottom, axis y line=none, ytick=\empty, grid=none,
      xtick={-10,-5,0,5,10}, xticklabels={−10,−5,0,+5,+10\,\%},
      xlabel={Fitted \(\tau\) vs \(R \cdot C\)},
      xlabel style={yshift=2pt}]
    \draw[chartText, line width=0.5pt, dash pattern=on 2pt off 1.5pt]
      (axis cs:-9.44,0.55) -- (axis cs:-9.44,4.75);
    \draw[chartText, line width=0.5pt] (axis cs:0,0.55) -- (axis cs:0,4.75);
    \node[note, anchor=base, yshift=3pt] at (axis cs:-9.44,4.75) {adjusted};
    \node[note, anchor=base, yshift=3pt] at (axis cs:0,4.75) {\(R \cdot C\)};
    \coordinate (prtop) at (axis cs:-12,4.75);
    \pairrow{1}{4}
    \pairrow{2}{3}
    \pairrow{3}{2}
    \pairrow{4}{1}
  \end{axis}
  \node[note, anchor=south west, inner sep=0pt] at ([yshift=14pt]prtop)
    {\chgmark\ charging\qquad\dismark\ discharging};
  \foreach \c/\rc in {1/47, 2/100, 3/147, 4/31.97} {
    % Tips at the bar's edges (half its 1.6 pt width off the line)
    \fill[traceInk] (prc\c) ++(0,-0.8pt) -- ++(2.5pt,-4.5pt) -- ++(-5pt,0) -- cycle;
    \fill[traceInk] (prd\c) ++(0,0.8pt) -- ++(2.5pt,4.5pt) -- ++(-5pt,0) -- cycle;
    \node[prhead] at ([yshift=10pt]prow\c)
      {\textbf{Circuit \c}\enspace\chgmark\,\prsec{c}{\c}{\rc}\enspace\dismark\,\prsec{d}{\c}{\rc}};
    \prdiff{\c}{\rc}{10pt}
    \prgap{\c}{1.5pt}   % 8.5 pt leading
  }
\end{tikzpicture}
